\documentclass[12pt]{article}



\usepackage{sbc-template}

\usepackage{graphicx,url}

\usepackage[utf8]{inputenc}

\usepackage[brazil]{babel}

\usepackage{array}

\usepackage{booktabs}

\usepackage{float}

\usepackage{enumitem}

\usepackage{amsmath}



\sloppy



\newcommand{\Nmod}{\ensuremath{N\_{\mathrm{mod}}}}

\newcommand{\NLOC}{\ensuremath{N\_{\mathrm{LOC}}}}

\newcommand{\CCN}{\ensuremath{CCN}}



\title{Priorização de Inspeção de Dívida Técnica por Hotspots Longitudinais em Software Público Brasileiro}



\author{Alan Mathias\inst{1}, Domini Acco\inst{1}, Guilherme Tombini\inst{1}, Jemison Santos\inst{1}}



\address{Universidade Estadual de Maringá (UEM)\\\\

Av. Colombo, 5790 -- Zona 07 -- 87020-900 -- Maringá -- PR -- Brasil\\\\

\email{pg2026009962\@uem.br}\\\\

\email{pg20261017594\@uem.br}\\\\

\email{pg20261018223\@uem.br}\\\\

\email{pg2026007135\@uem.br}}



\begin{document}

\maketitle



\begin{abstract}

Technical debt cannot be reliably diagnosed by a single code metric, which makes exhaustive inspection costly in large systems. This paper investigates whether change recurrence and structural complexity can be combined to prioritize files for technical-debt inspection in Brazilian public software. We conduct a longitudinal multi-case study with five public systems and annual snapshots between 2020 and 2026. For each source file, we extract modification frequency, lines of code, and cyclomatic complexity. Priority is represented by the continuous score $H=\min(R_N,R_C)$, which combines the relative positions of modification frequency and complexity without defining hotspots through a fixed percentage threshold. We analyze ranking stability over time and compare 25 high-priority files with 25 size- and history-aware controls through contextual inspection. Explicit contextual evidence was found in 28\% of high-priority files and 12\% of controls, but the paired comparison was inconclusive ($p=0.289$). A complementary paired static-analysis comparison found higher issue density in 9 hotspots and 11 controls, with 5 ties ($p=0.824$). The results support longitudinal hotspots as a search-space reduction mechanism, not as an automatic diagnosis of technical debt.

\end{abstract}



\begin{resumo}

Dívida técnica não pode ser diagnosticada de forma confiável por uma única métrica de código, o que torna custosa a inspeção exaustiva de sistemas de maior porte. Este trabalho investiga se recorrência de mudança e complexidade estrutural podem ser combinadas para priorizar arquivos para inspeção de dívida técnica em software público brasileiro. Foi realizado um estudo longitudinal multi-caso com cinco sistemas públicos e snapshots anuais entre 2020 e 2026. Para cada arquivo de código-fonte, foram extraídas frequência de modificação, linhas de código e complexidade ciclomática. A prioridade é representada pelo escore contínuo $H=\min(R_N,R_C)$, que combina as posições relativas de frequência de modificação e complexidade sem definir hotspots por um percentual fixo. Em seguida, avalia-se a estabilidade do ranking no tempo e comparam-se 25 arquivos de alta prioridade com 25 controles pareados por tamanho e histórico. Evidência contextual explícita foi encontrada em 28\% dos arquivos prioritários e em 12\% dos controles, embora a comparação pareada tenha sido inconclusiva ($p=0{,}289$). Em uma análise estática complementar pareada, 9 hotspots apresentaram maior densidade de ocorrências, contra 11 controles e 5 empates ($p=0{,}824$). Os resultados sustentam hotspots longitudinais como mecanismo de redução do espaço de busca, e não como diagnóstico automático de dívida técnica.

\end{resumo}



\section{Introdução}



Dívida técnica descreve compromissos técnicos que podem atender objetivos de curto prazo, mas aumentar custos de manutenção e evolução posteriormente. Estudos de mapeamento mostram que o fenômeno possui diferentes tipos, indicadores e formas de gerenciamento; por isso, uma métrica isolada não deve ser interpretada como diagnóstico automático de dívida técnica \cite{li2015mapping,alves2016identification}.



Uma alternativa é usar informações já presentes no processo de evolução do software para decidir \emph{onde inspecionar primeiro}. O histórico de mudanças contém informação relevante sobre manutenção e incidência de problemas \cite{graves2000changehistory,moser2008changemetrics}, enquanto a complexidade ciclomática descreve a estrutura do fluxo de controle \cite{mccabe1976complexity}. Essas dimensões podem ser combinadas para priorização sem assumir que alta mudança ou alta complexidade, isoladamente, representem dívida técnica.



Neste estudo, \emph{hotspot} significa uma região do código que merece inspeção prioritária por combinar elevada recorrência de modificação e elevada complexidade estrutural relativa. O objetivo não é criar um \`\`escore de dívida técnica'', mas reduzir o espaço de busca para análise posterior. Essa distinção é coerente com a literatura de Self-Admitted Technical Debt (SATD), que utiliza comentários e outros artefatos registrados por desenvolvedores como evidência contextual \cite{potdar2014satd,li2023satd}, e com estudos que combinam histórico, estrutura e contexto para investigar dívida técnica ao longo da evolução do software \cite{sutoyo2025tracing}.



O contexto empírico é o software público brasileiro. A experiência brasileira de software público está associada ao compartilhamento de soluções e à modernização da administração pública \cite{alves2009spb,pereira2019spb}. A Lei n. 14.129/2021 também reforça a relação entre governo digital, modernização e eficiência na prestação de serviços públicos \cite{leigovernodigital}. Além da relevância institucional, repositórios públicos permitem auditar e reproduzir a análise.



O objetivo geral é investigar se hotspots longitudinais de mudança e complexidade constituem uma estratégia útil para priorizar a inspeção de possíveis evidências de dívida técnica em softwares públicos brasileiros. O estudo é organizado em três questões de pesquisa:



\begin{description}[leftmargin=1.3cm,style=nextline]

\item[RQ1.] Como recorrência de modificação e complexidade estrutural são combinadas para produzir uma prioridade de inspeção?

\item[RQ2.] Quão estável é essa prioridade ao longo da evolução dos sistemas?

\item[RQ3.] Arquivos de alta prioridade apresentam mais evidências contextuais compatíveis com dívida técnica do que arquivos comparáveis de baixa prioridade?

\end{description}



\section{Trabalhos relacionados}



\subsection{Dívida técnica e evidência contextual}



\cite{li2015mapping} e \cite{alves2016identification} mostram que dívida técnica é um conceito amplo e que seus indicadores precisam ser interpretados no contexto do sistema. Essa literatura fundamenta a principal cautela deste trabalho: mudança e complexidade são sinais para priorização, não prova da existência de dívida.



A literatura de SATD fornece uma ponte entre métricas quantitativas e evidência contextual. \cite{potdar2014satd} mostraram que comentários de código podem registrar trabalho incompleto, soluções temporárias e necessidades de retrabalho. \cite{li2023satd} ampliaram essa perspectiva para comentários, commits, pull requests e issues. Assim, a inspeção de arquivos priorizados pode procurar evidências explícitas sem assumir que palavras-chave isoladas sejam suficientes para classificar dívida técnica.



\subsection{Mudança, complexidade e evolução}



Métricas derivadas do histórico de mudanças são recorrentes em estudos de qualidade de software. \cite{graves2000changehistory} mostraram que o histórico de mudanças contém informação relevante sobre incidência de falhas, e \cite{moser2008changemetrics} observaram valor informativo em métricas de processo. Neste trabalho, a frequência de modificação é utilizada como medida de recorrência de manutenção.



A \CCN{} deriva da medida proposta por \cite{mccabe1976complexity} e representa uma dimensão estrutural do código. A combinação de histórico e estrutura é coerente com estudos longitudinais recentes sobre dívida técnica, que também ressaltam a necessidade de evidência contextual para interpretação \cite{sutoyo2025tracing}. \NLOC{} é utilizada apenas para tornar a comparação entre arquivos prioritários e controles mais equilibrada em termos de tamanho.



\section{Metodologia}



\subsection{Desenho e corpus}



A pesquisa é um estudo empírico longitudinal multi-caso, baseado em mineração de repositórios Git e inspeção contextual. A descoberta dos projetos partiu de fontes públicas, incluindo o diretório GitHub Government Community \cite{githubgovernment}, e foi ampliada para organizações públicas brasileiras. O processo resultou em 77 itens inspecionados, 45 softwares catalogados, 21 candidatos elegíveis ao núcleo analítico e cinco sistemas selecionados.



A seleção final foi intencional, buscando diversidade de porte, domínio, maturidade, tecnologia e contexto institucional. Esse tipo de escolha é compatível com estudos de caso e com recomendações de amostragem em engenharia de software, desde que os critérios de seleção e os limites de generalização sejam explicitados \cite{runeson2009casestudy,baltes2022sampling}. A unidade quantitativa de análise é o arquivo de código-fonte \emph{first-party}; dependências de terceiros, artefatos gerados, arquivos de build e testes são excluídos conforme as regras de cada projeto.



\begin{table}[H]

\centering

\caption{Corpus longitudinal analisado.}

\label{tab:corpus}

\small

\begin{tabular}{lrrrrl}

\toprule

\textbf{Sistema} & \textbf{Período} & \textbf{Anos} & \textbf{Arq. únicos} & \textbf{Obs. arq.-ano} & \textbf{Linguagem} \\\\

\midrule

Novo SGP & 2020--2026 & 7 & 9.637 & 50.840 & C\\# \\\\

SAPL & 2020--2026 & 7 & 239 & 1.411 & Python \\\\

SIGA & 2020--2025 & 6 & 2.115 & 10.895 & Java \\\\

SIGI & 2020--2026 & 7 & 299 & 1.516 & Python \\\\

Painel e-SUS APS & 2024--2025 & 2 & 599 & 906 & Python/JS/TS \\\\

\bottomrule

\end{tabular}

\end{table}



\subsection{Snapshots, métricas e escore de prioridade}



Foram utilizados snapshots anuais entre 2020 e 2026, conforme a disponibilidade de histórico de cada sistema. A mineração usa Git e PyDriller \cite{pydriller}; tamanho e complexidade são obtidos com Lizard \cite{lizard}. Para cada arquivo são consideradas três medidas: \Nmod{} (número de modificações na janela), \NLOC{} (tamanho no snapshot) e \CCN{} (complexidade ciclomática). O escore de prioridade utiliza apenas \Nmod{} e \CCN{}.



Para evitar um corte arbitrário como \`\`top 20\\%'', \Nmod{} e \CCN{} são convertidos em posições relativas dentro de cada combinação sistema--linguagem--snapshot. Sejam $R_N(f,t)$ e $R_C(f,t)$ os postos relativos, reescalados no intervalo $[0,1]$. Define-se:



\begin{equation}

H(f,t)=\min\left(R_N(f,t),R_C(f,t)\right).

\end{equation}



O mínimo implementa uma regra conjuntiva: prioridade elevada exige que o arquivo esteja simultaneamente bem posicionado em recorrência de mudança e complexidade. Para resumir a prioridade ao longo do histórico de cada arquivo, utiliza-se a mediana dos valores observados:



\begin{equation}

\widetilde{H}(f)=\mathrm{mediana}\_{t}\\,H(f,t).

\end{equation}



Assim, $H$ e $\widetilde H$ são escores de priorização, não classificadores de dívida técnica.



\subsection{Procedimento para as questões de pesquisa}



\textbf{RQ1 -- priorização.} O escore $H$ é calculado em cada snapshot e $\widetilde H$ resume a prioridade longitudinal. Os arquivos são ordenados por $\widetilde H$, de modo que valores maiores indiquem maior combinação de recorrência de mudança e complexidade relativa. Não há um percentual fixo que defina o constructo de hotspot.



\textbf{RQ2 -- estabilidade longitudinal.} A estabilidade do ranking $H$ entre snapshots consecutivos é estimada pela correlação de postos de Spearman, considerando apenas arquivos presentes nos dois momentos. O objetivo é verificar se a lista de prioridades pode ser tratada como estática ou precisa ser recalculada ao longo do tempo.



\textbf{RQ3 -- evidência contextual.} Para a inspeção foram formados 25 pares, cinco por sistema, com um arquivo da região de alta prioridade e um controle de menor prioridade. O pareamento foi realizado dentro do mesmo sistema e linguagem, considerando proximidade em tamanho (\NLOC{}) e quantidade de snapshots observados. Os cortes usados apenas para compor essa amostra são operacionais e não definem hotspot. O código do último snapshot observado foi submetido a triagem de \texttt{TODO}/\texttt{FIXME}; ocorrências positivas foram lidas no contexto da implementação. O desfecho principal é binário: presença ou ausência de evidência explícita de manutenção pendente, necessidade de refatoração, risco conhecido ou implementação incompleta. As proporções pareadas são comparadas por teste binomial exato sobre pares discordantes.



\subsection{Análise estática complementar}

Além da inspeção contextual empregada na RQ3, foi realizada uma análise estática complementar sobre a mesma amostra pareada de 25 arquivos de alta prioridade e 25 controles. Essa etapa teve como objetivo verificar se arquivos priorizados pelo escore longitudinal também apresentavam maior concentração de ocorrências detectadas por ferramentas automatizadas de qualidade e manutenção de código.

A análise preservou integralmente os pares formados para a RQ3. Portanto, não houve nova seleção de arquivos com base nos resultados das ferramentas estáticas. Para cada par, hotspot e controle pertencem ao mesmo sistema e linguagem e foram previamente pareados considerando tamanho e quantidade de snapshots observados.

As ferramentas foram selecionadas de acordo com a linguagem dos arquivos. Para Python, foram utilizados Pylint e Radon; para Java, PMD; e, para C\#, os analisadores disponíveis no ecossistema .NET/Roslyn. O Pylint foi utilizado para identificar ocorrências classificadas como erros, avisos, refatorações e convenções, enquanto o Radon forneceu métricas de complexidade ciclomática e índice de manutenibilidade. O PMD foi empregado para identificar ocorrências relacionadas a design, boas práticas, possíveis erros, desempenho e complexidade. Para C\#, a inspeção foi realizada durante a restauração e compilação dos respectivos projetos com os analisadores .NET habilitados.

Como arquivos maiores tendem a apresentar maior quantidade absoluta de ocorrências, foi calculada a densidade de achados por mil linhas de código:

\begin{equation}
D_{\mathrm{issues}}(f)=\frac{\mathrm{issues}(f)}{\mathrm{LOC}(f)}\times 1000.
\end{equation}

A comparação foi realizada exclusivamente dentro de cada par. Para cada par, registrou-se se a densidade era maior no hotspot, maior no controle ou igual entre os dois arquivos. Como PMD, Pylint e Roslyn possuem conjuntos de regras e escalas distintas, valores absolutos não foram comparados diretamente entre linguagens ou ferramentas.

Alguns arquivos pertencentes à amostra original já não estavam presentes no HEAD atual dos respectivos repositórios. Para preservar a amostra previamente definida e evitar substituições \emph{post hoc}, nesses casos foi recuperada a revisão Git mais recente em que o caminho exato do arquivo ainda existia. O arquivo foi analisado nessa revisão, e o identificador da revisão utilizada foi registrado no pacote de replicação.

A direção dos pares discordantes foi avaliada por teste exato bilateral de sinais. Essa análise é exploratória e verifica apenas se uma das duas classes de arquivos tende a apresentar maior densidade de ocorrências dentro dos pares, não constituindo teste direto da presença de dívida técnica.


\section{Resultados}



\subsection{RQ1 -- prioridade de inspeção}



A aplicação de $H$ gera um ranking contínuo em cada snapshot e, por meio de $\widetilde H$, uma prioridade longitudinal por arquivo. Isso permite ordenar o código sem declarar que uma fronteira percentual específica separa \`\`hotspots'' de \`\`não hotspots''. A Tabela\~\ref{tab:rq1} apresenta, para tornar o resultado concreto, o arquivo com maior $\widetilde H$ entre os cinco arquivos de alta prioridade selecionados para inspeção em cada sistema.



\begin{table}[H]

\centering

\caption{Exemplos de arquivos de alta prioridade utilizados na inspeção contextual.}

\label{tab:rq1}

\small

\begin{tabular}{lp{7.2cm}r}

\toprule

\textbf{Sistema} & \textbf{Arquivo} & \textbf{$\widetilde H$} \\\\

\midrule

Novo SGP & \texttt{ObterNotasFrequenciaUseCase.cs} & 0,999 \\\\

SAPL & \texttt{sessao/views.py} & 0,980 \\\\

SIGA & \texttt{ExMovimentacaoController.java} & 0,997 \\\\

SIGI & \texttt{convenios/models.py} & 0,964 \\\\

Painel e-SUS APS & \texttt{nominal\\\_list\\\_adapter.py} & 0,994 \\\\

\bottomrule

\end{tabular}

\end{table}



A RQ1, portanto, não procura demonstrar que esses arquivos possuem dívida técnica. Ela estabelece apenas uma ordem de inspeção baseada na interseção contínua entre recorrência de mudança e complexidade estrutural.



\subsection{RQ2 -- estabilidade longitudinal da prioridade}



A estabilidade do ranking variou entre os sistemas (Tabela\~\ref{tab:rq2}). O Painel e-SUS APS apresentou a maior correlação mediana, 0,695, mas possui apenas uma transição anual comparável. SAPL e SIGI apresentaram estabilidade intermediária, enquanto Novo SGP e SIGA ficaram próximos de 0,40.



\begin{table}[H]

\centering

\caption{RQ2 -- estabilidade de $H$ entre snapshots consecutivos.}

\label{tab:rq2}

\begin{tabular}{lrr}

\toprule

\textbf{Sistema} & \textbf{Spearman mediano} & \textbf{Transições válidas} \\\\

\midrule

Novo SGP & 0,403 & 6 \\\\

SAPL & 0,539 & 6 \\\\

SIGA & 0,394 & 4 \\\\

SIGI & 0,574 & 6 \\\\

Painel e-SUS APS & 0,695 & 1 \\\\

\bottomrule

\end{tabular}

\end{table}



Os resultados mostram alguma persistência, mas não estabilidade suficiente para tratar a priorização como uma lista permanente. Em uso prático, o ranking precisa ser recalculado à medida que o sistema evolui.



\subsection{RQ3 -- evidência contextual nos arquivos priorizados}



A inspeção encontrou evidência contextual explícita em 7 dos 25 arquivos de alta prioridade (28\\%) e em 3 dos 25 controles (12\\%). Houve seis pares com evidência apenas no arquivo prioritário e dois apenas no controle; o teste exato sobre os pares discordantes resultou em $p=0{,}289$. A diferença observada segue a direção esperada, mas a amostra não permite concluir que arquivos prioritários apresentem evidência com frequência estatisticamente distinta dos controles.



\begin{table}[H]

\centering

\caption{RQ3 -- presença de evidência contextual nos 25 pares inspecionados.}

\label{tab:rq3}

\small

\begin{tabular}{lrr}

\toprule

\textbf{Sistema} & \textbf{Alta prioridade} & \textbf{Controle} \\\\

\midrule

Novo SGP & 0/5 & 0/5 \\\\

SAPL & 4/5 & 2/5 \\\\

SIGA & 3/5 & 1/5 \\\\

SIGI & 0/5 & 0/5 \\\\

Painel e-SUS APS & 0/5 & 0/5 \\\\

\midrule

Total & 7/25 (28\\%) & 3/25 (12\\%) \\\\

\bottomrule

\end{tabular}

\end{table}



O sinal ficou concentrado em SAPL e SIGA. Entre os casos encontrados aparecem comentários que identificam \texttt{HACK}, necessidade de refatoração, risco conhecido de erro, tratamento de exceção pendente e métodos incompletos ou \emph{stubs}. A RQ3 fornece, portanto, evidência direcional de utilidade para inspeção, mas não permite afirmar que arquivos prioritários \`\`contêm dívida técnica''.



\subsection{Análise estática complementar}

A análise estática foi concluída para os 25 pares da amostra, totalizando 50 arquivos. Em nove pares, o arquivo de alta prioridade apresentou maior densidade de ocorrências; em onze, a maior densidade foi observada no controle; e cinco pares resultaram em empate. Considerando os 20 pares discordantes, o teste exato bilateral de sinais resultou em $p=0{,}824$.

\begin{table}[H]
\centering
\caption{Comparação pareada da densidade de ocorrências da análise estática.}
\label{tab:static-paired}
\small
\begin{tabular}{lrrr}
\toprule
\textbf{Sistema} & \textbf{Hotspot maior} & \textbf{Controle maior} & \textbf{Empate} \\
\midrule
Novo SGP & 0 & 0 & 5 \\
SAPL & 1 & 4 & 0 \\
SIGA & 1 & 4 & 0 \\
SIGI & 3 & 2 & 0 \\
Painel e-SUS APS & 4 & 1 & 0 \\
\midrule
Total & 9 & 11 & 5 \\
\bottomrule
\end{tabular}
\end{table}

Os resultados não indicam uma tendência global de maior densidade de ocorrências estáticas nos arquivos de alta prioridade. A distribuição entre hotspots e controles foi aproximadamente equilibrada, e a comparação pareada não apresentou evidência de diferença sistemática entre os grupos.

Esse resultado não significa ausência de problemas estruturais nos arquivos priorizados. Em casos individuais, foram encontrados sinais expressivos. No SIGA, por exemplo, \texttt{ExMovimentacaoController.java} foi classificado pelo PMD como possível \emph{God Class}, apresentou acoplamento entre objetos igual a 102 e complexidade ciclomática total de 893, com máximo individual igual a 53.

Entretanto, arquivos controle do mesmo sistema também apresentaram evidências estruturais relevantes. \texttt{SrSolicitacao.java}, por exemplo, foi igualmente classificado como possível \emph{God Class}, com acoplamento igual a 94 e complexidade ciclomática total de 895. Outro controle, \texttt{SrRelPrazo.java}, apresentou um método com complexidade cognitiva de 488 e complexidade ciclomática de 135.

Assim, a análise estática evidencia que problemas estruturais podem ocorrer tanto em arquivos de alta prioridade quanto em controles. O resultado reforça que a quantidade ou densidade de ocorrências produzidas por analisadores estáticos não equivale diretamente à quantidade de dívida técnica de um arquivo.


\section{Discussão}

A principal contribuição deste estudo é operacional: utilizar recorrência de modificação e complexidade estrutural para produzir uma fila contínua de inspeção. O escore $H$ não pretende diagnosticar dívida técnica, mas reduzir o espaço de busca e indicar arquivos que merecem investigação mais detalhada.

Os resultados da RQ2 mostram que essa prioridade é dinâmica. As correlações entre snapshots consecutivos foram moderadas e variaram entre os sistemas, indicando que a lista de arquivos prioritários não deve ser considerada permanente. Em um contexto de uso prático, o ranking precisa acompanhar a evolução do software.

Na RQ3, arquivos de alta prioridade apresentaram evidência contextual explícita em 28\% dos casos, contra 12\% dos controles. A diferença foi observada na direção esperada, mas permaneceu inconclusiva na comparação pareada ($p=0{,}289$). Assim, a evidência contextual sugere possível utilidade do ranking como filtro de investigação, mas não demonstra que arquivos de alta prioridade contenham dívida técnica com frequência superior aos controles.

A análise estática complementar oferece uma perspectiva independente e ajuda a delimitar essa interpretação. Ao contrário da diferença observada na inspeção contextual, a densidade de ocorrências estáticas não favoreceu os hotspots: nove pares apresentaram maior densidade no arquivo prioritário, onze no controle e cinco resultaram em empate. O teste exato de sinais não indicou diferença entre as direções ($p=0{,}824$).

Esse resultado é relevante porque mostra que a prioridade longitudinal e a quantidade de ocorrências produzidas por analisadores estáticos representam dimensões distintas. Um arquivo pode combinar elevada recorrência de mudança e elevada complexidade relativa sem necessariamente disparar mais regras de uma ferramenta estática do que um arquivo controle. Da mesma forma, controles podem apresentar problemas estruturais expressivos.

Os casos do SIGA ilustram essa distinção. O hotspot \texttt{ExMovimentacaoController.java} apresentou sinais fortes de complexidade e acoplamento, incluindo indicação de possível \emph{God Class}. Entretanto, arquivos do grupo controle também apresentaram padrões semelhantes ou, em alguns casos, métricas estruturais ainda mais elevadas. Portanto, um caso individual com forte convergência entre prioridade e análise estática não pode ser generalizado para toda a amostra.

A triangulação dos resultados produz, assim, uma interpretação consistente. A análise longitudinal responde à pergunta sobre \emph{onde inspecionar primeiro}; a evidência contextual verifica se existem registros explícitos compatíveis com trabalho pendente, risco ou necessidade de refatoração; e a análise estática acrescenta indicadores automatizados de qualidade e estrutura. Nenhuma dessas fontes, isoladamente, deve ser tratada como diagnóstico definitivo de dívida técnica.

Do ponto de vista prático, o método pode ser utilizado como mecanismo de triagem em sistemas de grande porte. Em vez de inspecionar indiscriminadamente todo o repositório, uma equipe pode utilizar o ranking longitudinal para selecionar candidatos e, posteriormente, combinar análise humana, comentários, histórico de desenvolvimento, ferramentas estáticas e outros artefatos para decidir se existe dívida técnica relevante e qual sua natureza.


\section{Ameaças à validade}

\paragraph{Construto.} \Nmod{}, \CCN{} e $H$ não medem dívida técnica diretamente. Neste trabalho, essas variáveis são utilizadas exclusivamente para priorização. A evidência contextual baseada em comentários também é incompleta: ausência de marcadores não implica ausência de dívida, e a presença de \texttt{TODO}, \texttt{FIXME} ou termos semelhantes exige interpretação no contexto da implementação.

A análise estática também não constitui medida direta de dívida técnica. PMD, Pylint, Radon e Roslyn possuem regras, categorias e escalas distintas. Por esse motivo, suas contagens absolutas não são utilizadas como escala comum entre sistemas ou linguagens. A densidade de ocorrências é comparada apenas dentro de cada par formado no mesmo sistema e linguagem.

\paragraph{Validade interna.} Arquivos maiores podem apresentar mais ocorrências simplesmente por conterem mais código. Para reduzir esse efeito, os controles foram pareados considerando \NLOC{} e quantidade de snapshots, e a análise estática adicional utiliza densidade de ocorrências por mil linhas de código.

Alguns arquivos da amostra já não estavam disponíveis no HEAD atual dos repositórios. Nesses casos, foi recuperada a revisão Git mais recente em que o caminho exato ainda existia, preservando a identidade dos arquivos originalmente amostrados. Essa estratégia evita substituição \emph{post hoc}, embora introduza diferenças temporais entre alguns arquivos analisados.

Em um dos arquivos C\# recuperados historicamente, a restauração das dependências foi concluída, mas a compilação da revisão histórica retornou erro. Como esse caso resultou em empate de zero ocorrências e não integra os pares discordantes utilizados no teste de sinais, sua exclusão não altera a direção da conclusão estatística. Ainda assim, esse caso deve ser interpretado com cautela.

\paragraph{Validade externa.} Os cinco sistemas constituem amostra intencional e não representam estatisticamente todo o universo de software público brasileiro. A generalização pretendida é analítica e se restringe ao comportamento do procedimento em sistemas com diferentes domínios, tecnologias, portes e históricos de evolução \cite{baltes2022sampling}.

\paragraph{Confiabilidade.} Repositórios, identificadores de revisão, filtros, scripts, arquivos de configuração, amostra pareada e resultados intermediários foram preservados no pacote de replicação. Os resultados da análise estática registram também a ferramenta utilizada, o modo de resolução do arquivo e a revisão de origem, permitindo distinguir arquivos avaliados no HEAD daqueles recuperados historicamente.


\section{Conclusão}

Este trabalho investigou se a combinação longitudinal entre recorrência de modificação e complexidade estrutural pode ser utilizada para priorizar a inspeção de possíveis evidências de dívida técnica em software público brasileiro.

O escore proposto permitiu produzir uma prioridade contínua de inspeção sem definir hotspots por um percentual fixo. A análise longitudinal mostrou, contudo, que essa prioridade não é estática: sua estabilidade variou entre os sistemas, indicando a necessidade de recalcular o ranking à medida que o software evolui.

Na inspeção contextual, evidências explícitas foram encontradas em 28\% dos arquivos de alta prioridade e em 12\% dos controles. Embora a diferença tenha ocorrido na direção esperada, a comparação pareada permaneceu inconclusiva ($p=0{,}289$).

A análise estática complementar dos mesmos 25 pares não mostrou maior densidade de ocorrências nos hotspots. Dos 25 pares, nove favoreceram os arquivos prioritários, onze favoreceram os controles e cinco apresentaram empate. Entre os pares discordantes, o teste exato de sinais resultou em $p=0{,}824$, não indicando diferença sistemática entre os grupos.

Esses resultados reforçam a distinção central deste estudo: \textbf{hotspots longitudinais não constituem detectores de dívida técnica}. Sua utilidade está em reduzir o espaço de busca e estabelecer uma ordem de inspeção baseada na combinação entre pressão evolutiva e complexidade estrutural. A existência, a natureza e a relevância da dívida técnica dependem de evidências adicionais.

A triangulação entre histórico de mudanças, estrutura do código, evidência contextual e análise estática mostrou que essas fontes capturam dimensões distintas. Alguns arquivos de alta prioridade apresentam sinais estruturais e contextuais fortes, enquanto outros não; da mesma forma, arquivos controle podem conter problemas estáticos relevantes.

Assim, a principal conclusão do estudo é que recorrência de modificação e complexidade estrutural podem ser utilizadas para responder à pergunta \emph{``onde inspecionar primeiro?''}, mas não à pergunta \emph{``onde existe dívida técnica?''}. Como continuidade, a investigação pode incorporar commits, issues, pull requests e decisões arquiteturais, além de ampliar a amostra para avaliar se determinados tipos de dívida ou categorias de problemas apresentam maior associação com a prioridade longitudinal.


\bibliographystyle{sbc}

\bibliography{references}

\end{document}
